\subsection{Geographic and Structured Validation}
Cross-validation design should reflect the structure of the observations and the intended use of a predictive model. Classical studies of cross-validation and bootstrap evaluation provide the statistical context for evaluating predictive errors \citep{kohavi1995,arlot2010}. Roberts et al.~\citep{roberts2017} examine validation strategies for data with spatial, temporal, hierarchical, and related dependence, emphasizing that randomly held-out observations can provide a different estimate of generalization than grouped or spatially separated samples. In health-related prediction, Subbaswamy and Saria~\citep{subbaswamy2020} discuss dataset shift and the challenge of maintaining model performance when deployment conditions differ from those in development. These studies motivate our leave-one-area-out experiment, but neither establishes that any particular state must have higher error than another.

\subsection{Medicare Payment Standardization}
The CMS DME geography-and-service dataset supplies the observed payment and standardized payment fields used in this study \citep{cmsdme2023}. O'Donnell et al.~\citep{odonnell2013} examine the rationale for standardizing Medicare payments when making geographic comparisons: observed reimbursements can include differences attributable to payment policy as well as utilization and service patterns. Their study uses a different Medicare population and is not a prior experiment on DME code-level prediction. Consequently, we compare the two CMS target fields descriptively and predictively without assuming that standardized DME payments eliminate all regional differences.

\subsection{Interpretable Baselines and Evaluation}
Prediction evaluation also requires attention to the measures reported and the consequences of model-selection bias. Steyerberg and colleagues review a framework for assessment of prediction models \citep{steyerberg2010}, while Cawley and Talbot examine over-fitting during model selection \citep{cawley2010}. Varma and Simon~\citep{varma2006} document bias when cross-validation is also used for choosing a model. These works motivate retaining transparent prespecified baselines without asserting that they are optimal.

Before fitting more flexible algorithms, a group-mean predictor can establish how much information is conveyed by a categorical identifier. We first estimate mean payments for HCPCS codes, then estimate means for joint HCPCS/rental groups with prespecified fallbacks for unobserved categories. These baselines provide a transparent comparison in which errors can be attributed to observed and unseen codes rather than to opaque model behavior. Computation uses Python and the open-source numerical and machine-learning ecosystem, including scikit-learn \citep{pedregosa2011}. The contribution of this study is the combined evaluation protocol and its observed results, not the invention of a new prediction algorithm.